\chapter*{Cara Membaca: Empat Jalur}
\addcontentsline{toc}{chapter}{Cara Membaca: Empat Jalur}
\markboth{Cara Membaca}{Cara Membaca}

Setiap bab punya bentuk yang sama, sehingga pembaca tahu apa yang diharapkan dan
boleh melompat tanpa kehilangan arah.

\begin{center}
\small
\begin{tabular}{@{}p{3.1cm}p{4.6cm}p{4.4cm}@{}}
\toprule
\textbf{Jalur} & \textbf{Yang dibaca} & \textbf{Hasil yang didapat} \\
\midrule
Pembaca muda (10-14) & kotak \emph{Untuk pembaca muda}, Ringkasan, satu contoh
  hitung & gagasan inti tanpa notasi; tahu apa yang sedang dihitung \\
Mahasiswa S1 & seluruh badan bab, Lab komputasi, Latihan; lewati kotak peneliti &
  definisi, bukti, kemampuan menghitung dan memeriksa \\
Dosen & kotak \emph{Catatan untuk dosen}, peta sumber, Latihan & rencana
  pertemuan, miskonsepsi yang harus diantisipasi, bahan papan tulis \\
Peneliti & kotak \emph{Untuk peneliti}, peta sumber, Lampiran~A dan~B & batas
  klaim, sumber primer, pertanyaan terbuka, notasi yang konsisten \\
\bottomrule
\end{tabular}
\end{center}

\medskip
Tanda \textsc{tingkat} di awal setiap bagian menunjukkan kesulitan bacaan:
\emph{dasar} dapat dibaca tanpa bab sebelumnya selain yang ditunjuk;
\emph{menengah} menuntut aljabar dan notasi buku ini;
\emph{lanjut} boleh dilewati pada pembacaan pertama tanpa memutus alur bab.

\medskip
\noindent\textbf{Enam cara belajar, enam pintu masuk.} Orang tidak belajar dengan cara
yang sama. Karena itu setiap bab memuat lebih dari satu rute menuju gagasan yang sama,
dan tidak ada rute yang wajib: pilih yang paling cepat membuatmu paham, lalu kembali ke
badan bab untuk melihat mengapa cara itu bekerja.

\begin{center}
\small
\begin{tabular}{@{}p{3.4cm}p{4.3cm}p{4.6cm}@{}}
\toprule
\textbf{Tipe} & \textbf{Cari ini di setiap bab} & \textbf{Cara memakainya} \\
\midrule
Visual & 2-3 gambar bernomor (di Daftar Gambar), tabel, dan warna kotak & baca
  gambar sebelum teksnya; gambar memuat angka yang sama dengan hitungan di teks \\[2pt]
Auditori & kotak \emph{Bacakan keras-keras} (indigo) & bacakan naskahnya keras-keras
  atau rekam suaramu; ulangi bagian yang berirama \\[2pt]
Kinestetik & kotak \emph{Kerjakan dengan tangan} (hijau) & lakukan dengan benda nyata:
  kertas, kartu, lantai, uang, atau badan sendiri \\[2pt]
Pemaham konsep & definisi, teorema, bukti, kotak \emph{Untuk peneliti} & jangan pakai
  rumus sebelum tahu asalnya; bukti adalah tempat asal itu ditulis \\[2pt]
Penghafal prosedural & kotak \emph{Resep langkah} (cokelat) dan Latihan & hafalkan
  urutan kartu resep, lalu terapkan pada angka berbeda sampai otomatis \\[2pt]
Pemecah masalah & kotak \emph{Tantangan} (ungu) dan butir latihan bertanda &
  kerjakan tantangan lebih dulu; petunjuknya bertingkat, jangan dibuka sekaligus \\
\bottomrule
\end{tabular}
\end{center}

Empat jalur di bawah ini mengatur \emph{berapa banyak} yang dibaca; enam tipe di atas
mengatur \emph{lewat pintu mana} masuk. Keduanya boleh dikombinasikan: seorang pembaca
visual yang mengejar jalur peneliti tetap membaca semua bukti, hanya saja ia mulai dari
gambarnya.

Tiga jalur belajar yang sudah dicoba, dari yang paling ringan:

\begin{enumerate}
  \item \textbf{Jalur rasa ingin tahu (8-12 minggu).} Baca kotak oranye, Ringkasan,
        dan satu contoh hitung dari setiap bab, berurutan dari Bab~1 sampai Bab~32.
        Tanpa latihan. Tujuannya satu: tahu peta medan.
  \item \textbf{Jalur mahasiswa (2 semester).} Semester pertama: Bab~1-10 dan
        Bab~17. Semester kedua: Bab~11-16 lalu 18-32. Itu dua puluh satu bab, lebih
        banyak dari satu semester biasa; cara memangkasnya tanpa memutus prasyarat
        adalah menunda Bab~31-32 (kriptografi dan privasi) ke mata kuliah pilihan,
        sehingga tersisa sembilan belas bab. Kerjakan Lab komputasi setiap bab;
        latihan minimal empat butir per bab. Urutan ini menjaga janji prasyarat:
        Bab~18 menuntut Bab~13 dan~17, jadi ketiganya tidak boleh berjauhan; Bab~21
        menuntut Bab~10 dan~20; Bab~23 menuntut Bab~12.
  \item \textbf{Jalur peneliti keamanan ML (9 minggu).} Sembilan belas bab, bukan
        tujuh: Bab~3, 4, 9, 10, 12, 13, 18, 19, 20, 22, 23, 25, lalu 26-30, ditutup
        Bab~31 dan~32. Bab~13, 18, 22, dan~23 ikut karena masing-masing adalah
        prasyarat Bab~20 (uji hipotesis), Bab~29 (adversarial), dan Bab~25 (jaringan
        saraf); Bab~31 dan~32 melengkapi sisi kriptografi dan privasi yang dituntut
        judul buku. Membaca jalur ini dalam enam minggu berarti melewati prasyarat
        dan kehilangan alasan mengapa klaimnya berlaku.
\end{enumerate}

Simbol dipakai konsisten di seluruh buku: $\mathbb{N}$, $\mathbb{Z}$, $\mathbb{Q}$,
$\mathbb{R}$ untuk himpunan bilangan; $\mathbf{x}$ untuk vektor dan $x_i$ untuk
komponennya; $\|\mathbf{x}\|$ untuk norma Euklides; $\nabla f$ untuk gradien;
$\varepsilon$ dan $\delta$ selalu bilangan positif kecil; $\mathbb{E}[X]$ untuk
nilai harapan; $\mathcal{L}(V,W)$ untuk ruang transformasi linear; $\top$ untuk
transpos. Lampiran~A memuat daftar lengkapnya.
